\def\bagrutTocEntry{\phantomsection\addcontentsline{toc}{subsection}{\texorpdfstring{سؤال $16$ امتحان $\mathrm{899381}$ سنة $2022$}{سؤال 16 امتحان 899381 سنة 2022}}\global\let\bagrutTocEntry\relax}%
\includepdf[pages=-,pagecommand={\thispagestyle{empty}\bagrutTocEntry}]{../../../bagrut_questions/oop/oop_2022_899381_16.pdf}


\ifwithsols
\begin{boxSolution}[\texorpdfstring{سؤال $16$ امتحان $\mathrm{899381}$ سنة $2022$ - فرع أ}{سؤال 16 امتحان 899381 سنة 2022 - فرع أ}]
\textbf{أ.}
\begin{enumerate}[label=\textbf{(\arabic*)}, itemsep=0.8em]
    \item المصفوفة \textenglish{arr} تحوي قمصانًا وبنطلونات، وبما أنّ الفئتين \textenglish{Shirt} و \textenglish{Pants} ترثان من الفئة \textenglish{Clothing}، يكون نمط المصفوفة \textenglish{Clothing[]}:
\begin{boxCode}
\begin{english}
\begin{minted}{csharp}
public class Cart
{
    private string name;
    private Clothing[] arr;
    private int current;
}
\end{minted}
\end{english}
\end{boxCode}

    \item المخطّط الهرميّ للعلاقة بين الفئات (الفئتان \textenglish{Shirt} و \textenglish{Pants} ترثان من الفئة \textenglish{Clothing}، والفئة \textenglish{Cart} تحوي مصفوفة من نمط \textenglish{Clothing}):
\begin{center}
\begin{english}
\begin{tikzpicture}[
    classbox/.style={rectangle, draw, thick, minimum width=2.4cm, minimum height=0.75cm, font=\ttfamily\bfseries},
    inherit/.style={-{Triangle[open, length=2.6mm, width=2.6mm]}, thick, shorten >=1pt},
    contain/.style={-{Diamond[length=3.2mm, width=2.2mm]}, thick}
]
    \node[classbox] (clothing) at (0, 0) {Clothing};
    \node[classbox] (cart) at (-4.2, 0) {Cart};
    \node[classbox] (shirt) at (-1.8, -1.5) {Shirt};
    \node[classbox] (pants) at (1.8, -1.5) {Pants};

    \draw[contain] (clothing.west) -- (cart.east);
    \draw[inherit] (shirt.north) -- ([xshift=-0.4cm]clothing.south);
    \draw[inherit] (pants.north) -- ([xshift=0.4cm]clothing.south);
\end{tikzpicture}
\end{english}
\end{center}
\end{enumerate}
\end{boxSolution}

\filbreak

\begin{boxSolution}[\texorpdfstring{سؤال $16$ امتحان $\mathrm{899381}$ سنة $2022$ - فرع ب (1) و(2)}{سؤال 16 امتحان 899381 سنة 2022 - فرع ب (1) و(2)}]
\textbf{ب.}
\begin{enumerate}[label=\textbf{(\arabic*)}, itemsep=0.8em]
    \item المخطّط الهرميّ المعدّل بعد إضافة الفئتين \textenglish{Product} و \textenglish{Book}:
\begin{center}
\begin{english}
\begin{tikzpicture}[
    classbox/.style={rectangle, draw, thick, minimum width=2.4cm, minimum height=0.75cm, font=\ttfamily\bfseries},
    inherit/.style={-{Triangle[open, length=2.6mm, width=2.6mm]}, thick, shorten >=1pt},
    contain/.style={-{Diamond[length=3.2mm, width=2.2mm]}, thick}
]
    \node[classbox] (product) at (0, 0) {Product};
    \node[classbox] (cart) at (-4.2, 0) {Cart};
    \node[classbox] (clothing) at (-1.8, -1.5) {Clothing};
    \node[classbox] (book) at (1.8, -1.5) {Book};
    \node[classbox] (shirt) at (-3.3, -3.0) {Shirt};
    \node[classbox] (pants) at (-0.3, -3.0) {Pants};

    \draw[contain] (product.west) -- (cart.east);
    \draw[inherit] (clothing.north) -- ([xshift=-0.4cm]product.south);
    \draw[inherit] (book.north) -- ([xshift=0.4cm]product.south);
    \draw[inherit] (shirt.north) -- ([xshift=-0.4cm]clothing.south);
    \draw[inherit] (pants.north) -- ([xshift=0.4cm]clothing.south);
\end{tikzpicture}
\end{english}
\end{center}

    \item الصفتان المشتركتان لجميع المنتجات (الملابس والكتب) هما رقم التشخيص \textenglish{id} والسعر \textenglish{price}، وتُنقلان إلى الفئة الأساسية \textenglish{Product}:
\begin{boxCode}
\begin{english}
\begin{minted}{csharp}
public class Product
{
    private int id;
    private double price;
}
\end{minted}
\end{english}
\end{boxCode}
\end{enumerate}
\end{boxSolution}

\filbreak

\begin{boxSolution}[\texorpdfstring{سؤال $16$ امتحان $\mathrm{899381}$ سنة $2022$ - فرع ب (3)}{سؤال 16 امتحان 899381 سنة 2022 - فرع ب (3)}]
\begin{enumerate}[label=\textbf{(\arabic*)}, start=3]
    \item الفئات التي طرأ فيها تغيير هي \textenglish{Clothing} و \textenglish{Book} (ترثان من \textenglish{Product} وتُحذف منهما الصفتان \textenglish{id} و \textenglish{price})، والفئة \textenglish{Cart} (يتغيّر نمط المصفوفة \textenglish{arr} إلى \textenglish{Product[]}):
\begin{boxCode}
\begin{english}
\begin{minted}{csharp}
public class Clothing : Product
{
    private string fabric;
    private string color;
}

public class Book : Product
{
    private string bookName;
    private string author;
}

public class Cart
{
    private string name;
    private Product[] arr;
    private int current;
}
\end{minted}
\end{english}
\end{boxCode}
\end{enumerate}
\end{boxSolution}

\filbreak

\begin{boxSolution}[\texorpdfstring{سؤال $16$ امتحان $\mathrm{899381}$ سنة $2022$ - فرع جـ}{سؤال 16 امتحان 899381 سنة 2022 - فرع جـ}]
\textbf{جـ.}
\begin{itemize}[itemsep=0.4em]
    \item في الفئة \textenglish{Product}: نطبّق العملية \textenglish{GetDiscountPrice} لتعيد السعر بعد تخفيض $10\%$ (أي ضرب السعر بـ $0.9$).
    \item في الفئة \textenglish{Clothing}: نعيد تعريف العملية (\textenglish{override}) بحيث تستدعي عملية الفئة الأساسية وتطرح من السعر $10$ شواكل إضافية (وبذلك ترث الفئتان \textenglish{Shirt} و \textenglish{Pants} هذا التطبيق تلقائيًّا).
\end{itemize}
\begin{boxCode}
\begin{english}
\begin{minted}{csharp}
// In class Product:
public virtual double GetDiscountPrice()
{
    return this.price * 0.9;
}

// In class Clothing:
public override double GetDiscountPrice()
{
    return base.GetDiscountPrice() - 10;
}
\end{minted}
\end{english}
\end{boxCode}
\end{boxSolution}
\fi
\filbreak
